\section{Introduction}
A decision interface can appear reliable under aggregate accuracy while changing its answers when the same candidate set is presented in a different order. This discrepancy matters when an application routes a request, triggers a review, or selects an action from a finite schema. The application consumes the selected action, rather than the model's average score over unrelated examples. Evaluation must therefore ask both whether a decision matches its reference and whether equivalent presentations preserve that decision.

System1Bench studies interfaces that map a supplied state and typed questions to structured answers. The motivating systems include Laya and Jev, alongside language models constrained to the same answer spaces \citep{layamodel,jevdocs}. The term ``System 1'' identifies this product/interface category. Our unit of comparison is a \emph{model plus adapter on a fixed downstream decision}, not an entire deployed workflow. Whole-system success also depends on retrieval, tools, human review and action execution, none of which we measure.

The motivating empirical result is concrete. On the same 100 Banking77 control states, reversing candidate presentation changes 48\% of Laya English predictions while increasing reference accuracy by nine percentage points. Laya Multilingual changes 38\% of predictions with zero net accuracy change. A leaderboard displaying only the two marginal accuracies would omit most of the behavioral change. The comparison is not a universal degradation claim: it identifies sensitivity that has no fixed direction in aggregate correctness.

Our approach combines a cross-domain source atlas with matched perturbation diagnostics. Sources differ in application domain, decision function, reference provenance, output type, candidate cardinality, language, context length and the number of questions sharing a state. We preserve these distinctions instead of averaging unlike references. Paired tests compare the same semantic unit across encodings, and same-order repeats distinguish changed presentation from ordinary execution variability. Separate operational experiments measure resident-model completion and hosted calls under their respective observable boundaries.

This paper makes three contributions. First, it maps 15 sources to non-exclusive application domains and decision functions while retaining reference provenance and sample identity. Second, it separates net correctness, correctness turnover and raw prediction discordance under paired interventions, including a frozen policy diagnostic and an orthogonal answer-code control. A subsequent held-out wording experiment tests and qualifies a mechanism suggested by exploratory inspection. Third, it supplies a reproducible quality-and-efficiency record with frozen local timing controls and a separately pinned hosted API track. The historical analyses are exploratory: their outputs existed before the diagnostic questions were finalized; the two generated follow-ups were frozen before their respective inference runs.

The central implication is methodological. Stable accuracy, stable decisions and efficient execution are distinct properties. A practical benchmark should reveal which of them a system satisfies on a particular workload.
